\documentclass[reqno]{amsart}

\usepackage{lmodern}
\usepackage[T1]{fontenc}
\usepackage[utf8]{inputenc}
\usepackage[english]{babel}
\usepackage{microtype}
\usepackage{amsmath,amssymb,amsthm,mathrsfs,mathtools,mathdots}
\usepackage{xcolor}
\usepackage{booktabs,enumerate}
\usepackage[colorlinks=true,linkcolor=blue,citecolor=blue,urlcolor=blue]{hyperref}

\numberwithin{equation}{section}

\newtheorem{theorem}{Theorem}[section]
\newtheorem{proposition}[theorem]{Proposition}
\newtheorem{lemma}[theorem]{Lemma}
\newtheorem{corollary}[theorem]{Corollary}
\newtheorem{conjecture}[theorem]{Conjecture}

\theoremstyle{definition}
\newtheorem{definition}[theorem]{Definition}
\newtheorem{example}[theorem]{Example}
\newtheorem{remark}[theorem]{Remark}

\newcommand{\defin}[1]{%
\relax\ifmmode%
\textcolor{blue}{#1}%
\else\textcolor{blue}{\emph{#1}}%
\fi%
}

\newcommand{\setN}{\mathbb{N}}
\newcommand{\setR}{\mathbb{R}}
\newcommand{\poly}{\setR[t]}
\newcommand{\coeff}[2]{[#1]#2}
\newcommand{\had}{\odot}
\newcommand{\opI}{\mathcal{I}}
\newcommand{\pf}{\mathrm{PF}}
\DeclareMathOperator{\bl}{bl}
\DeclareMathOperator{\roots}{roots}

\title[Super-recurrence Eulerian polynomials]{Real-rootedness of a
super-recurrence Eulerian family}
\author{Per Alexandersson}
\date{\today}

\subjclass[2020]{Primary 05A15; Secondary 26C10, 05A20}
\keywords{Eulerian polynomial, real-rooted polynomial, interlacing,
Hadamard product, P{\'o}lya frequency sequence}

\begin{document}

\begin{abstract}
We prove real-rootedness for a family of Eulerian-type polynomials introduced
by U.~Shankar.  For an integer $l\geq 1$, the coefficients are defined by
\[
  E_{1,0}^{(l)}=1,\qquad
  E_{n,k}^{(l)}
  =(k+1)^lE_{n-1,k}^{(l)}
  +(n-k)^lE_{n-1,k-1}^{(l)} .
\]
We show that the row polynomial
$E_n^{(l)}(t)=\sum_{k=0}^{n-1}E_{n,k}^{(l)}t^k$ has only real
nonpositive zeros.  In fact, we prove the stronger consecutive interlacing
relation
\[
  E_n^{(l)}(t)\ll E_{n+1}^{(l)}(t)
\]
for all $l\geq 1$ and $n\geq 1$.  The proof is based on a Hadamard
normalization of the rows, the Garloff--Wagner theorem on Hadamard products,
and a small kernel interlacing argument.
\end{abstract}

\maketitle

\section{Introduction}
\label{sec:introduction}

Eulerian polynomials are a central source of real-rooted and interlacing
phenomena in enumerative combinatorics.  One way to see this is through the
classical recurrence
\[
  A_{n,k}=(k+1)A_{n-1,k}+(n-k)A_{n-1,k-1}.
\]
Shankar~\cite{Shankar2025SuperRecurrence} considered a broad family of
triangular arrays obtained by replacing the linear factors in such recurrences
by powers.  The Eulerian specialization is the following.

\begin{definition}
\label{def:superEulerianRows}
Fix $l\geq 1$.  We define the \defin{super-recurrence Eulerian
coefficients} $E_{n,k}^{(l)}$ by
\[
  E_{1,0}^{(l)}=1
\]
and, for $n\geq 2$,
\begin{equation}
\label{eq:superEulerianRecurrence}
  E_{n,k}^{(l)}
  =(k+1)^lE_{n-1,k}^{(l)}
  +(n-k)^lE_{n-1,k-1}^{(l)} ,
\end{equation}
where coefficients outside the range $0\leq k\leq n-1$ are interpreted as
zero.  The \defin{row polynomial} is
\[
  E_n^{(l)}(t)
  \coloneqq
  \sum_{k=0}^{n-1}E_{n,k}^{(l)}t^k .
\]
\end{definition}

For $l=1$ this is the usual Eulerian recurrence.  For example, when $l=2$ the
first rows are
\[
  E_1^{(2)}(t)=1,\qquad
  E_2^{(2)}(t)=1+t,\qquad
  E_3^{(2)}(t)=1+8t+t^2,
\]
and
\[
  E_4^{(2)}(t)=1+41t+41t^2+t^3 .
\]
Shankar proved palindromicity and log-concavity properties for these arrays
and recorded real-rootedness evidence.  The main result of this note settles
the real-rootedness question in the Eulerian specialization, and gives a
stronger interlacing statement.

\begin{theorem}
\label{thm:main}
Let $l\geq 1$.
\begin{enumerate}[(a)]
\item For every $n\geq 1$, the polynomial $E_n^{(l)}(t)$ has only real
  nonpositive zeros.
\item For every $n\geq 1$,
  \[
    E_n^{(l)}(t)\ll E_{n+1}^{(l)}(t).
  \]
\end{enumerate}
\end{theorem}

Here $f\ll g$ denotes the usual weak proper-position relation for real-rooted
polynomials, with the orientation fixed in Section~\ref{sec:tools}.  The proof
also gives a refinement for a natural family of cumulative block-leader
prefixes.  We state it now and define the prefixes in
Section~\ref{sec:prefixes}.

\begin{theorem}
\label{thm:prefixChainIntro}
For $l\geq 1$, $n\geq 1$, and $1<m\leq n$, the cumulative block-leader
prefixes satisfy
\[
  P_{n,m-1}^{(l)}(t)\ll P_{n,m}^{(l)}(t).
\]
\end{theorem}

Our proof has two main ideas.  The first is a normalization.  If $d=n-1$, put
\[
  Q_n^{(l)}(t)
  \coloneqq
  \sum_{i=0}^{d}
  \frac{E_{n,i}^{(l)}}{\binom di^l}t^i .
\]
Then
\[
  E_n^{(l)}(t)=B_{d,l}(t)\had Q_n^{(l)}(t),
  \qquad
  B_{d,l}(t)=\sum_{i=0}^{d}\binom di^l t^i ,
\]
where $\had$ denotes coefficientwise Hadamard product.  The kernel
$B_{d,l}$ is the $l$-fold Hadamard power of $(1+t)^d$, while the normalized
rows $Q_n^{(l)}$ satisfy a recurrence built from the two operators
$N-\theta$ and $\theta+1$, where $\theta=t\,d/dt$.  These operators preserve
the cone of polynomials with nonnegative coefficients and only nonpositive
real zeros.  This proves part~(a) of Theorem~\ref{thm:main}.

The second idea is that the same normalization is compatible with interlacing.
We use the two-pair form of the Garloff--Wagner Hadamard theorem
\cite[Thm.~4]{GarloffWagner1996}.  A small pair of explicit kernels,
\[
  V_d(t)=(1+t)^d,\qquad
  U_d(t)=\sum_{i=0}^{d}(i+1)\binom di t^i
        =(1+t)^{d-1}(1+(d+1)t),
\]
contains all the needed root-list information.  Hadamard powers of the
relations among $V_d$, $U_d$, and the shifted reciprocal of $U_d$ then yield
the consecutive interlacing in part~(b).

\section{Real-rootedness tools}
\label{sec:tools}

We collect the standard facts about real-rooted polynomials that are used in
the proof.  For background on P{\'o}lya frequency sequences, real-rootedness,
compatibility, and interlacing, see
Aissen--Edrei--Schoenberg--Whitney~\cite{AissenEdreiSchoenbergWhitney1951},
Brenti~\cite{Brenti1989}, Chudnovsky--Seymour~\cite{ChudnovskySeymour2007},
and Br{\"a}nd{\'e}n~\cite{Branden2015}.

\begin{definition}
\label{def:pf}
A polynomial $f\in\poly$ is a \defin{PF polynomial} if all its coefficients
are nonnegative and all its zeros are real and nonpositive.  We use the
convention that the zero polynomial is allowed when it appears as a limiting
or degenerate case.
\end{definition}

If $f,g\in\poly$ have only real zeros, we write $f\ll g$ if the zeros of $f$
weakly interlace or weakly alternate to the left of the zeros of $g$, in the
following orientation.  If $\deg g=\deg f+1$, the roots can be listed as
\[
  \beta_1\leq \alpha_1\leq \beta_2\leq
  \dotsb\leq \alpha_{\deg f}\leq \beta_{\deg g},
\]
where the $\alpha_i$ are roots of $f$ and the $\beta_i$ are roots of $g$.  If
$\deg f=\deg g$, the roots can be listed as
\[
  \alpha_1\leq \beta_1\leq \alpha_2\leq
  \dotsb\leq \alpha_{\deg f}\leq \beta_{\deg g}.
\]
We use the closed weak form, so repeated roots and zero polynomials are
handled by the usual limiting convention.

For $f(t)=\sum_i a_it^i$ and $g(t)=\sum_i b_it^i$, define their
\defin{Hadamard product}
\[
  (f\had g)(t)\coloneqq \sum_i a_ib_i t^i ,
\]
where coefficients are extended by zero outside the degree range.  We also
write
\[
  \opI_Df(t)\coloneqq t^D f(1/t)
\]
for the \defin{shifted reciprocal}.

\begin{theorem}[Schur--P{\'o}lya--Wagner; Garloff--Wagner]
\label{thm:hadamardTools}
Let $f,g,p,q\in\poly$ be PF polynomials.
\begin{enumerate}[(a)]
\item If $f$ and $p$ are PF, then $f\had p$ is PF.
\item If $f\ll g$ and $p\ll q$, then
  \[
    f\had p\ll g\had q.
  \]
\end{enumerate}
\end{theorem}

Part~(a) is the classical Schur--P{\'o}lya--Wagner preservation theorem, in the
nonnegative-coefficient convention used here.  Part~(b) is the two-pair
proper-position form of Garloff--Wagner~\cite[Thm.~4]{GarloffWagner1996};
see also the earlier work of Schur and P{\'o}lya~\cite{SchurPolya1914}.

\begin{theorem}[Fixed-left cone property]
\label{thm:fixedLeftCone}
Let $p,q,r\in\poly$ be PF polynomials.  If $p\ll q$ and $p\ll r$, then
\[
  p\ll aq+br
\]
for all $a,b\geq 0$.
\end{theorem}

This is the standard convexity of a common interlacing cone.  It is equivalent
to the two-generator compatibility criterion of Obreschkoff type; in the form
used here it is also contained in the addition lemmas used by
Garloff--Wagner~\cite[Lem.~6]{GarloffWagner1996}.

We shall use two elementary differential operators.  Let
\[
  \theta f(t)\coloneqq t f'(t).
\]

\begin{lemma}
\label{lem:operatorsPreservePF}
Let $f\in\poly$ be PF.
\begin{enumerate}[(a)]
\item The polynomial $(\theta+1)f=(tf)'$ is PF.
\item If $\deg f\leq N$, then $(N-\theta)f$ is PF.
\end{enumerate}
\end{lemma}

\begin{proof}
The coefficients are plainly nonnegative in both cases, in the second case
because the coefficient of $t^i$ in $(N-\theta)f$ is $(N-i)\coeff{t^i}f$.

For the roots, part~(a) follows from Rolle's theorem applied to $tf$.  Indeed,
the polynomial $tf$ has only real nonpositive zeros, and its derivative is
$(\theta+1)f$.

For part~(b), remove the zero roots of $f$.  Thus write
\[
  f(t)=t^s h(t),\qquad
  h(t)=c\prod_{j=1}^{m}(t+a_j),
\]
where $c>0$ and $a_j>0$.  Since $s+m=\deg f\leq N$, we have
\[
  (N-\theta)f
  =
  t^s\bigl((N-s-\theta)h\bigr).
\]
Moreover,
\[
  (N-s-\theta)h
  =
  (N-s-m)h+\sum_{j=1}^{m}a_j\,\frac{h}{t+a_j}.
\]
The polynomials $h/(t+a_j)$ weakly interlace $h$, and all coefficients in
this linear combination are nonnegative.  By
Theorem~\ref{thm:fixedLeftCone}, the right-hand side is again real-rooted
with only nonpositive zeros.  Multiplying by $t^s$ only appends zero roots.
This completes the proof.
\end{proof}

\begin{lemma}
\label{lem:thetaPlusOnePreservesInterlacing}
The operator $\theta+1$ preserves weak proper position on the PF cone.  In
particular, if $f$ and $g$ are PF and $f\ll g$, then
\[
  (\theta+1)f\ll(\theta+1)g.
\]
Consequently, every power $(\theta+1)^r$ has the same property.
\end{lemma}

\begin{proof}
The compatibility characterization of proper position says that
$f\ll g$ is equivalent, in this nonnegative PF setting, to real-rootedness of
all nonnegative linear combinations $af+bg$.  Applying
Lemma~\ref{lem:operatorsPreservePF}(a) to each such combination gives that
\[
  a(\theta+1)f+b(\theta+1)g
\]
is PF for all $a,b\geq 0$.  The same characterization gives the desired
proper-position relation.  Iterating proves the final assertion.
\end{proof}

\section{The normalized rows}
\label{sec:normalizedRows}

The proof of Theorem~\ref{thm:main} starts with a normalization which removes
the binomial part of the coefficients.

\begin{definition}
\label{def:normalizedRows}
For $d=n-1$, define
\[
  Q_n^{(l)}(t)
  \coloneqq
  \sum_{i=0}^{d}
  \frac{E_{n,i}^{(l)}}{\binom di^l}t^i .
\]
We also define
\[
  B_{d,l}(t)\coloneqq \sum_{i=0}^{d}\binom di^l t^i .
\]
\end{definition}

\begin{lemma}[Normalized transfer identity]
\label{lem:normalizedTransfer}
For $l\geq 1$ and $n\geq 2$,
\[
  Q_n^{(l)}(t)
  =
  (1+t)
  \left(
    \frac{(\theta+1)(n-1-\theta)}{n-1}
  \right)^l
  Q_{n-1}^{(l)}(t).
\]
\end{lemma}

\begin{proof}
Put $N=n-1$ and write
\[
  Q_{n-1}^{(l)}(t)=\sum_{k=0}^{N-1}q_kt^k .
\]
The operator
\[
  \frac{(\theta+1)(N-\theta)}{N}
\]
is diagonal in the monomial basis: it sends $t^k$ to
\[
  \frac{(k+1)(N-k)}{N}t^k .
\]
Thus, after applying the $l$th power and multiplying by $1+t$, the coefficient
of $t^k$ is
\begin{equation}
\label{eq:normalizedTransferCoeff}
  \left(\frac{(k+1)(N-k)}{N}\right)^l q_k
  +
  \left(\frac{k(N+1-k)}{N}\right)^l q_{k-1},
\end{equation}
where $q_{-1}=q_N=0$.

Since
\[
  q_k=\frac{E_{n-1,k}^{(l)}}{\binom{N-1}{k}^l},
\]
the first summand in \eqref{eq:normalizedTransferCoeff} is
\[
  \frac{(k+1)^lE_{n-1,k}^{(l)}}{\binom Nk^l},
\]
using
\[
  \binom Nk=\frac{N}{N-k}\binom{N-1}{k}.
\]
Similarly, the second summand is
\[
  \frac{(N+1-k)^lE_{n-1,k-1}^{(l)}}{\binom Nk^l},
\]
using
\[
  \binom Nk=\frac{N}{k}\binom{N-1}{k-1}.
\]
Since $N+1=n$, the sum is
\[
  \frac{
    (k+1)^lE_{n-1,k}^{(l)}
    +(n-k)^lE_{n-1,k-1}^{(l)}
  }{\binom{n-1}{k}^l}
  =
  \frac{E_{n,k}^{(l)}}{\binom{n-1}{k}^l}.
\]
This is exactly the coefficient of $t^k$ in $Q_n^{(l)}$.
\end{proof}

\begin{lemma}
\label{lem:normalizedRowsPF}
For all $l\geq 1$ and $n\geq 1$, the polynomial $Q_n^{(l)}$ is PF.
\end{lemma}

\begin{proof}
We proceed by induction on $n$.  The base case is $Q_1^{(l)}=1$.  Assume
that $Q_{n-1}^{(l)}$ is PF.  Its degree is at most $n-2$.  Therefore
Lemma~\ref{lem:operatorsPreservePF}(b), with $N=n-1$, shows that
$(n-1-\theta)Q_{n-1}^{(l)}$ is PF.  Lemma~\ref{lem:operatorsPreservePF}(a)
then shows that applying $\theta+1$ preserves the PF property.  The scalar
factor $1/(n-1)$ is positive, and multiplication by $1+t$ only adds the root
$-1$.  Iterating this $l$ times in
Lemma~\ref{lem:normalizedTransfer} proves that $Q_n^{(l)}$ is PF.
\end{proof}

\begin{lemma}
\label{lem:rowHadamardFactorization}
For $d=n-1$,
\[
  E_n^{(l)}=B_{d,l}\had Q_n^{(l)}.
\]
Consequently, $E_n^{(l)}$ is PF for all $l\geq 1$ and $n\geq 1$.
\end{lemma}

\begin{proof}
The coefficient of $t^i$ in $B_{d,l}\had Q_n^{(l)}$ is
\[
  \binom di^l\frac{E_{n,i}^{(l)}}{\binom di^l}
  =
  E_{n,i}^{(l)}.
\]
This proves the factorization.

Moreover, $B_{d,l}$ is the $l$-fold Hadamard power of $(1+t)^d$, since
\[
  (1+t)^d=\sum_{i=0}^{d}\binom di t^i.
\]
The polynomial $(1+t)^d$ is PF, and so
Theorem~\ref{thm:hadamardTools}(a) implies that $B_{d,l}$ is PF.  By
Lemma~\ref{lem:normalizedRowsPF}, the other factor $Q_n^{(l)}$ is PF.
Applying Theorem~\ref{thm:hadamardTools}(a) once more proves that
$E_n^{(l)}$ is PF.
\end{proof}

This proves Theorem~\ref{thm:main}(a).

\section{Palindromicity and the Hadamard kernels}
\label{sec:kernels}

We now prepare the interlacing proof.  We first record the palindromicity,
which is also proved in Shankar's work~\cite{Shankar2025SuperRecurrence}.

\begin{lemma}
\label{lem:palindromic}
For $0\leq k\leq n-1$,
\[
  E_{n,k}^{(l)}=E_{n,n-1-k}^{(l)}.
\]
Consequently, both $E_n^{(l)}$ and $Q_n^{(l)}$ are palindromic of degree
$n-1$.
\end{lemma}

\begin{proof}
We proceed by induction on $n$.  The claim is immediate for $n=1$.  Assume it
for row $n-1$.  Applying the recurrence to the index $n-1-k$ gives
\[
  E_{n,n-1-k}^{(l)}
  =
  (n-k)^lE_{n-1,n-1-k}^{(l)}
  +(k+1)^lE_{n-1,n-2-k}^{(l)}.
\]
By the induction hypothesis in row $n-1$,
\[
  E_{n-1,n-1-k}^{(l)}=E_{n-1,k-1}^{(l)}
  \quad\text{and}\quad
  E_{n-1,n-2-k}^{(l)}=E_{n-1,k}^{(l)},
\]
with the usual zero convention at the boundary.  Hence
\[
  E_{n,n-1-k}^{(l)}
  =
  (n-k)^lE_{n-1,k-1}^{(l)}
  +(k+1)^lE_{n-1,k}^{(l)}
  =
  E_{n,k}^{(l)}.
\]
The assertion for $Q_n^{(l)}$ follows because
$\binom{n-1}{k}=\binom{n-1}{n-1-k}$.
\end{proof}

Define the second Hadamard kernel by
\[
  K_{d,l}(t)\coloneqq
  \sum_{i=0}^{d}(i+1)^l\binom di^l t^i .
\]
We shall compare $B_{d,l}$, $K_{d,l}$, and the shifted reciprocal
$\opI_{d+1}K_{d,l}$.

\begin{lemma}
\label{lem:kernelInterlacings}
For all $d\geq 0$ and $l\geq 1$,
\[
  B_{d,l}\ll K_{d,l},
  \qquad
  B_{d,l}\ll \opI_{d+1}K_{d,l},
  \qquad
  K_{d,l}\ll \opI_{d+1}K_{d,l}.
\]
\end{lemma}

\begin{proof}
The case $d=0$ is immediate, so assume $d\geq 1$.  Put
\[
  V_d(t)=(1+t)^d
\]
and
\[
  U_d(t)=\sum_{i=0}^{d}(i+1)\binom di t^i.
\]
Since
\[
  U_d(t)=(\theta+1)V_d(t)=(t(1+t)^d)'
        =(1+t)^{d-1}(1+(d+1)t),
\]
the roots of $U_d$ are $-1$, with multiplicity $d-1$, and
$-1/(d+1)$.  The polynomial $\opI_{d+1}U_d$ has roots $0$,
$-(d+1)$, and $-1$ with multiplicity $d-1$.  Comparing these ordered root
lists gives
\[
  V_d\ll U_d,
  \qquad
  V_d\ll \opI_{d+1}U_d,
  \qquad
  U_d\ll \opI_{d+1}U_d.
\]

Now
\[
  B_{d,l}=V_d^{\had l},
  \qquad
  K_{d,l}=U_d^{\had l},
  \qquad
  \opI_{d+1}K_{d,l}=(\opI_{d+1}U_d)^{\had l}.
\]
The last identity follows coefficientwise from the definition of the shifted
reciprocal.  Repeated application of
Theorem~\ref{thm:hadamardTools}(b) to the three displayed interlacings gives
the result.
\end{proof}

\section{Consecutive interlacing}
\label{sec:consecutive}

We now prove the stronger part of Theorem~\ref{thm:main}.  The proof is a
direct decomposition of the two summands in the recurrence for
$E_{n+1}^{(l)}$.

\begin{lemma}
\label{lem:boundaryDecomposition}
Let $n\geq 1$, put $d=n-1$, and set
\[
  p(t)=E_n^{(l)}(t),
  \qquad
  A(t)=(\theta+1)^lp(t).
\]
Then
\[
  E_{n+1}^{(l)}(t)=A(t)+t^nA(1/t).
\]
\end{lemma}

\begin{proof}
The recurrence for row $n+1$ gives
\[
  E_{n+1}^{(l)}(t)
  =
  \bigl((\theta+1)^l+t(n-\theta)^l\bigr)p(t).
\]
The first summand is $A(t)$.  We show that the second summand is
$t^nA(1/t)$.  Write
\[
  p(t)=\sum_{i=0}^{d}a_i t^i .
\]
By Lemma~\ref{lem:palindromic}, $a_i=a_{d-i}$.  The coefficient of $t^j$ in
$t(n-\theta)^lp(t)$ is
\[
  (n-j+1)^l a_{j-1}
\]
for $1\leq j\leq n$, and it is zero for $j=0$.  On the other hand,
the coefficient of $t^j$ in $t^nA(1/t)$ is
\[
  (n-j+1)^l a_{n-j}
  =
  (n-j+1)^l a_{j-1},
\]
again for $1\leq j\leq n$.  Thus the two polynomials agree.
\end{proof}

\begin{lemma}
\label{lem:summandDecompositions}
With the notation of Lemma~\ref{lem:boundaryDecomposition},
\[
  p=B_{d,l}\had Q_n^{(l)},
  \qquad
  A=K_{d,l}\had Q_n^{(l)},
\]
and
\[
  t^nA(1/t)
  =
  (\opI_{d+1}K_{d,l})\had tQ_n^{(l)}.
\]
\end{lemma}

\begin{proof}
The first identity is Lemma~\ref{lem:rowHadamardFactorization}.  For the
second, the coefficient of $t^i$ in $A=(\theta+1)^lp$ is
\[
  (i+1)^l E_{n,i}^{(l)}
  =
  (i+1)^l\binom di^l \coeff{t^i}{Q_n^{(l)}},
\]
which is the coefficient of $t^i$ in $K_{d,l}\had Q_n^{(l)}$.

For the last identity, compare coefficients.  Since $n=d+1$, the coefficient
of $t^j$ in the right-hand side is
\[
  (d+2-j)^l\binom d{d+1-j}^l\coeff{t^{j-1}}{Q_n^{(l)}}.
\]
Using
\[
  \binom d{d+1-j}=\binom d{j-1}
\]
and palindromicity of $Q_n^{(l)}$, this is exactly the coefficient of
$t^j$ in $t^nA(1/t)$.
\end{proof}

\begin{proof}[Proof of Theorem~\ref{thm:main}(b)]
For $n=1$, the assertion is $1\ll 1+t$, which is immediate.  Assume
$n\geq 2$ and use the notation of Lemma~\ref{lem:boundaryDecomposition}.
By Lemma~\ref{lem:summandDecompositions},
\[
  p=B_{d,l}\had Q_n^{(l)},\qquad
  A=K_{d,l}\had Q_n^{(l)}.
\]
The normalized row $Q_n^{(l)}$ is PF by Lemma~\ref{lem:normalizedRowsPF}, so
\[
  Q_n^{(l)}\ll Q_n^{(l)}.
\]
Together with $B_{d,l}\ll K_{d,l}$ from
Lemma~\ref{lem:kernelInterlacings}, the two-pair Hadamard theorem gives
\[
  p\ll A.
\]

Similarly,
\[
  t^nA(1/t)
  =
  (\opI_{d+1}K_{d,l})\had tQ_n^{(l)}.
\]
Since $Q_n^{(l)}$ is PF, the root list immediately gives
\[
  Q_n^{(l)}\ll tQ_n^{(l)}.
\]
Combining this with
\[
  B_{d,l}\ll \opI_{d+1}K_{d,l}
\]
and applying Theorem~\ref{thm:hadamardTools}(b), we get
\[
  p\ll t^nA(1/t).
\]
The fixed-left cone property, Theorem~\ref{thm:fixedLeftCone}, now gives
\[
  p\ll A+t^nA(1/t).
\]
By Lemma~\ref{lem:boundaryDecomposition}, the right-hand side is
$E_{n+1}^{(l)}$.  This proves
\[
  E_n^{(l)}\ll E_{n+1}^{(l)},
\]
as desired.
\end{proof}

The same proof gives a reciprocal endpoint which will be useful for the prefix
refinement.

\begin{corollary}
\label{cor:boundaryEndpoint}
Let
\[
  A_n^{(l)}(t)\coloneqq (\theta+1)^lE_n^{(l)}(t).
\]
Then
\[
  A_n^{(l)}(t)\ll t^nA_n^{(l)}(1/t)
\]
and
\[
  A_n^{(l)}(t)\ll E_{n+1}^{(l)}(t).
\]
\end{corollary}

\begin{proof}
The first statement follows from the same Hadamard decomposition as above,
using
\[
  K_{d,l}\ll \opI_{d+1}K_{d,l}
  \qquad\text{and}\qquad
  Q_n^{(l)}\ll tQ_n^{(l)}.
\]
The second follows from the fixed-left cone property and
\[
  E_{n+1}^{(l)}(t)=A_n^{(l)}(t)+t^nA_n^{(l)}(1/t).
\]
This completes the proof.
\end{proof}

\section{Cumulative block-leader prefixes}
\label{sec:prefixes}

We finish with the prefix refinement.  We use a block-leader product formula
for the rows.  A \defin{block-leader set} is a subset $B\subseteq[n]$ with
$1\in B$.  For such a set, define
\[
  c_n(B)\coloneqq \prod_{i=1}^{n}w_i(B),
\]
where, if
\[
  r_i=|B\cap[i-1]|,
\]
then
\[
  w_i(B)=
  \begin{cases}
  i-r_i, & i\in B,\\
  r_i, & i\notin B.
  \end{cases}
\]
Equivalently, if $B=\{1=b_1<b_2<\dotsb<b_r\}$ and $b_{r+1}=n+1$, then
\[
  c_n(B)
  =
  \prod_{j=1}^{r}(b_j-j+1)
  \prod_{j=1}^{r}j^{b_{j+1}-b_j-1}.
\]

\begin{lemma}
\label{lem:blockLeaderFormula}
We have
\[
  E_n^{(l)}(t)
  =
  \sum_{\substack{B\subseteq[n]\\1\in B}}
  c_n(B)^l t^{|B|-1}.
\]
\end{lemma}

\begin{proof}
Let the right-hand side be $F_n^{(l)}(t)$.  We verify the same recurrence and
initial condition as in Definition~\ref{def:superEulerianRows}.  The case
$n=1$ is clear.

Fix $n\geq 2$ and split the block-leader sets according to whether $n$ belongs
to $B$.  If $n\notin B$ and $|B|=k+1$, then $B$ is a block-leader set in
$[n-1]$ and
\[
  c_n(B)=(k+1)c_{n-1}(B).
\]
These terms contribute
\[
  (k+1)^l E_{n-1,k}^{(l)}
\]
to the coefficient of $t^k$.

If $n\in B$, then $B'=B\setminus\{n\}$ is a block-leader set in $[n-1]$ with
$|B'|=k$.  In this case
\[
  c_n(B)=(n-k)c_{n-1}(B'),
\]
and the contribution to the coefficient of $t^k$ is
\[
  (n-k)^lE_{n-1,k-1}^{(l)}.
\]
Thus the coefficients of $F_n^{(l)}$ satisfy
\eqref{eq:superEulerianRecurrence}.  The claim follows.
\end{proof}

\begin{definition}
\label{def:prefixes}
For $1\leq m\leq n$, define the \defin{cumulative block-leader prefix}
\[
  P_{n,m}^{(l)}(t)
  \coloneqq
  \sum_{\substack{B\subseteq[n]\\1\in B,\ \max B\leq m}}
  c_n(B)^l t^{|B|-1}.
\]
\end{definition}

Note that $P_{n,n}^{(l)}=E_n^{(l)}$ by Lemma~\ref{lem:blockLeaderFormula}.

\begin{lemma}
\label{lem:prefixClosedForm}
For $1\leq m\leq n$,
\[
  P_{n,m}^{(l)}(t)
  =
  (\theta+1)^{l(n-m)}E_m^{(l)}(t).
\]
\end{lemma}

\begin{proof}
If $B\subseteq[m]$ and we view $B$ as a subset of $[n]$, then every index
$i>m$ is not in $B$ and has $r_i=|B|$.  Hence
\[
  c_n(B)=|B|^{n-m}c_m(B).
\]
Using Lemma~\ref{lem:blockLeaderFormula}, we get
\[
  P_{n,m}^{(l)}(t)
  =
  \sum_{\substack{B\subseteq[m]\\1\in B}}
  |B|^{l(n-m)}c_m(B)^l t^{|B|-1}.
\]
The operator $\theta+1$ multiplies the monomial $t^{|B|-1}$ by $|B|$.
Therefore $(\theta+1)^{l(n-m)}$ gives exactly the displayed factor
$|B|^{l(n-m)}$.  This proves the identity.
\end{proof}

\begin{proof}[Proof of Theorem~\ref{thm:prefixChainIntro}]
It is enough to prove that
\[
  P_{n,m}^{(l)}\ll P_{n,m+1}^{(l)}
  \qquad (m<n),
\]
and then replace $m$ by $m-1$.  By Corollary~\ref{cor:boundaryEndpoint},
\[
  (\theta+1)^lE_m^{(l)}\ll E_{m+1}^{(l)}.
\]
Apply the operator
\[
  (\theta+1)^{l(n-m-1)}
\]
to both sides.  This preserves weak proper position by
Lemma~\ref{lem:thetaPlusOnePreservesInterlacing}.  Using
Lemma~\ref{lem:prefixClosedForm}, the resulting relation is exactly
\[
  P_{n,m}^{(l)}\ll P_{n,m+1}^{(l)}.
\]
The claim follows.
\end{proof}

\begin{remark}
The normalized gamma-polynomial refinements suggested by computations are not
used in the proof above.  The Hadamard normalization and the small kernel
interlacing argument are sufficient for the real-rootedness and consecutive
interlacing statements.
\end{remark}

\bibliographystyle{alpha}
\bibliography{bibliography}

\end{document}
